\chapter{Bittajming}
\label{ch:tajming}

{\large\itshape Hur noder utan gemensam klocka läser samma bit samtidigt.\par}

\begin{chapterbox}{I det här kapitlet}
\begin{itemize}
  \item Varför en CAN-bit delas in i segment i stället för att bara räknas.
  \item Sampelpunkten, och varför den ligger sent i bitperioden.
  \item Synkronisering: hård vid SOF, mjuk vid varje flank därefter.
  \item Vad som händer när två noders oscillatorer inte går lika fort.
\end{itemize}
\end{chapterbox}

\section{Problemet}

Varje nod har sin egen klocka. Ingen ledning bär takten --- CAN har fyra signaler och ingen av dem
är en klocka.

Ändå måste alla noder vara överens om var varje bit börjar och slutar, och det på bitnivå:
arbitreringen i \kapref{ch:arbitrering} bygger på att två noder läser \emph{samma} bitintervall
samtidigt och drar slutsatser av vad de ser. Är de ur fas jämför de olika bitar, och resultatet
blir godtyckligt.

Lösningen har två delar: en gemensam \emph{definition} av var i bitperioden man läser, och en
mekanism som drar tillbaka noder som driver i väg.

\section{Bitperioden i segment}

En CAN-bit är inte ett odelat intervall utan delas in i fyra segment, mätta i
\textbf{tidskvanta} --- korta, lika långa steg som noden räknar med sin egen klocka.

\bookfigure{bit_timing}{Bitperiodens fyra segment. Sampelpunkten ligger mellan fas\-segment~1 och
2, sent i perioden.}{fig:timing}

\begin{booktable}{@{}llL@{}}
\thead{Segment} & \thead{Längd} & \thead{Vad det är till för} \\ \midrule
Synksegment & 1 tq & Här förväntas en flank ligga. \\
Propagationssegment & 1--8 tq & Täcker signalens gångtid ut på bussen och tillbaka. \\
Fassegment 1 & 1--8 tq & Före sampelpunkten; kan förlängas vid resynkronisering. \\
Fassegment 2 & 1--8 tq & Efter sampelpunkten; kan förkortas. \\
\end{booktable}

\textbf{Sampelpunkten} ligger mellan fassegment 1 och 2. Det är där, och bara där, biten läses.

\subsection{Varför sent i perioden}

Sampelpunkten läggs typiskt 70--80~\% in i bitperioden. Skälet är propagationssegmentet, och det
är samma skäl som gör att bussens längd begränsar hastigheten.

Under arbitreringen måste en sändande nod hinna läsa tillbaka vad \emph{den bortersta} noden
driver, innan biten är slut. Signalen ska alltså hinna ut och tillbaka. Propagationssegmentet är
den tiden, och sampelpunkten måste ligga efter den --- annars läser noden bussen innan den andra
nodens bit hunnit fram, och arbitreringen blir fel.

\begin{aside}
\note Det är därför bithastighet och kabellängd hänger ihop (\kapref{ch:buss}). Högre hastighet
ger kortare bitperiod, alltså kortare propagationssegment i absolut tid, alltså kortare tillåten
kabel. Sambandet är inte elektriskt utan logiskt: det handlar om vad arbitreringen kräver.
\end{aside}

\section{Synkronisering}

\subsection{Hård synkronisering}

Vid SOF --- den enda flank alla noder ser samtidigt när bussen varit ledig --- gör varje mottagare
en \textbf{hård synkronisering}: den nollställer sin biträknare, oavsett var den stod.

Det är en gång per frame, och det är den enda gången noderna får sätta klockan helt fritt.

\subsection{Resynkronisering}

Under framen räcker inte det. Två oscillatorer som skiljer sig en bråkdels procent driver isär
långsamt, och över en frame på 130 bitar blir det mätbart.

Därför resynkroniserar varje mottagare vid \emph{varje} recessiv-till-dominant-flank i strömmen.
Ligger flanken tidigare än väntat förkortas fassegment 2; ligger den senare förlängs fassegment 1.
Justeringen är begränsad till ett antal tidskvanta som kallas \emph{synchronization jump width}.

\textbf{Och nu faller bitstoppningen på plats.} Femregeln i \kapref{ch:stuffing} garanterar att
det aldrig går mer än fem bitperioder utan en flank, vilket sätter en övre gräns för hur långt
två noder kan hinna driva isär innan de dras tillbaka.

\begin{aside}
De två mekanismerna är alltså inte oberoende: bitstoppningen finns \emph{för} resynkroniseringens
skull. Ett protokoll utan stoppning skulle behöva antingen en klockledning eller betydligt
noggrannare oscillatorer.
\end{aside}

\section{Att räkna ut tajmingen}

För en konstruktion räknas två tal ut ur systemklockan och den önskade bithastigheten:

\[ \text{tick per bit} = \frac{f_{\text{system}}}{\text{bithastighet}} \]

Med en systemklocka på 50~MHz och 1~Mbit/s blir det 50 klockcykler per CAN-bit. Sampelpunkten vid
70~\% blir då tick 35.

\begin{cppcode}
// Derived, never written down as a literal: change the bit rate and both follow.
constexpr unsigned TicksPerBit{SystemClockHz / BitRateHz};
constexpr unsigned SampleTick{TicksPerBit * 7U / 10U};
\end{cppcode}

Att de \emph{räknas ut} och inte skrivs in är poängen. En konstruktion där \code{50} och
\code{35} står inskrivna på var sitt ställe är en konstruktion där den ena hinner ändras utan
den andra.

\section{Vad som händer när det går fel}

Två fel är värda att känna igen, eftersom de ser olika ut:

\textbf{Fel bithastighet.} Noderna är ur fas från första biten. Symptomet är att ingenting alls
fungerar --- inga frames tas emot, och den sändande noden rapporterar fel direkt. Det är det
snällare felet, för det syns omedelbart.

\textbf{Fel sampelpunkt.} Noderna kommer överens om de flesta bitar men inte alla. Symptomet är
sporadiska CRC- och stoppfel som blir värre med längre kabel och högre belastning. Det här är
det obehagliga felet, för det ser ut som en dålig kabel.

\section*{Övningar}

\exercise{Räkna tajmingen}{Räkning}
\label{ex:5:1}
En konstruktion har systemklockan 50~MHz.
\begin{enumerate}[a)]
  \item Hur många klockcykler går det på en CAN-bit vid 1~Mbit/s?
  \item Vid 500~kbit/s?
  \item Var ligger sampelpunkten i klockcykler räknat, vid 1~Mbit/s och 70~\%?
\end{enumerate}

\exercise{Sampelpunkten}{Resonemang}
\label{ex:5:2}
\begin{enumerate}[a)]
  \item Varför ligger sampelpunkten sent i bitperioden i stället för i mitten?
  \item Vad skulle gå fel under arbitreringen om den låg vid 20~\%?
  \item Vilken del av bitperioden skulle behöva växa om kabeln blev dubbelt så lång?
\end{enumerate}

\exercise{Stoppning och synk}{Resonemang}
\label{ex:5:3}
\begin{enumerate}[a)]
  \item Hur många bitperioder kan som mest passera utan en flank i en CAN-frame?
  \item Vilken mekanism garanterar det?
  \item Vad skulle hända med resynkroniseringen om bitstoppningen togs bort?
\end{enumerate}

\exercise{Två symptom}{Resonemang}
\label{ex:5:4}
Två noder ska prata med varandra.
\begin{enumerate}[a)]
  \item Nod A är inställd på 500~kbit/s och nod B på 1~Mbit/s. Vad ser ni?
  \item Båda kör 1~Mbit/s, men A samplar vid 50~\% och B vid 80~\%. Vad ser ni, och varför är
    det svårare att felsöka?
\end{enumerate}
